\documentclass[11pt,a4paper]{article}
\usepackage[a4paper,left=20mm,right=20mm,top=21mm,bottom=23mm]{geometry}
\usepackage[T2A]{fontenc}
\usepackage[utf8]{inputenc}
\usepackage[russian]{babel}
\usepackage{amsmath,amssymb,mathtools,booktabs,array,longtable}
\usepackage{xcolor,hyperref,tikz}
\hypersetup{colorlinks=true,linkcolor=blue!60!black,urlcolor=blue!60!black}
\setlength{\parindent}{0pt}
\setlength{\parskip}{4pt}
\setlength{\emergencystretch}{3em}
\renewcommand{\arraystretch}{1.17}
\newcommand{\Idea}{\textcolor{blue!60!black}{\textbf{Ключевая идея:}}}
\newcommand{\Res}{\textcolor{teal!60!black}{\textbf{Итог:}}}
\newcommand{\Heading}[2]{\par\medskip\textbf{#1}\par#2\par}
\begin{document}
\thispagestyle{empty}
\vspace*{0.30\textheight}
\begin{center}
{\LARGE\bfseries Ревью домашней работы\par}
\vspace{23mm}{\large Автор: Лёвин Артём Александрович\par}
\vspace{9mm}{\large 19 октября 2026 года\par}
\end{center}
\vfill
\begin{center}\begin{tikzpicture}[x=1cm,y=1cm]
\draw[blue!35!black,line width=.6pt] (0,0) .. controls (1.5,.32) and (3.5,-.38) .. (5.2,0) .. controls (6.9,.34) and (8.1,-.22) .. (9.5,.03);
\end{tikzpicture}\end{center}
\newpage
\section*{Сводная таблица ошибок}
\footnotesize
\begin{longtable}{@{}p{1.55cm}p{2.55cm}p{4.15cm}p{2.6cm}p{4.2cm}@{}}
\toprule
№ & Тема & Суть ошибки & Ваш ответ & Правильный ответ\\
\midrule\endfirsthead
\toprule № & Тема & Суть ошибки & Ваш ответ & Правильный ответ\\\midrule\endhead
\hyperref[task:1_7]{1-7} & Логарифмы и ОДЗ & Не проверены допустимые значения логарифмов & $0{,}5$ & $-0{,}5$ \\
\hyperref[task:1_9]{1-9} & Производная и экстремумы & Неверно посчитаны переходы знака производной & $3$ & $2$ \\
\hyperref[task:1_12]{1-12} & Корни и графики & Не завершено нахождение пересечения & Ответ не указан & $15$ \\
\hyperref[task:1_13]{1-13} & Кредит по годам & Вместо платежа за год найдена сумма всех выплат & $1{,}8264$ млн руб. & $356400$ руб. \\
\hyperref[task:1_14]{1-14} & Тригонометрия & Неверное разложение выражения на множители & Получены посторонние серии корней & $-3\pi+$\newline $\arctan(1/2)$\newline и\newline $-3\pi+$\newline $\arctan(11/2)$ \\
\hyperref[task:2_1]{2-1} & Площадь треугольника & Не использован синус угла между сторонами & $\frac12$ & $\frac32$ \\
\hyperref[task:2_5]{2-5} & Вероятность двух матчей & Не перечислены все подходящие сочетания результатов & На листе записаны $0{,}3$ и $0{,}5$ & $0{,}33$ \\
\hyperref[task:2_8]{2-8} & Логарифмы и дроби & Ошибка при сложении отрицательной дроби & $-3{,}5$ & $-2{,}5$ \\
\hyperref[task:2_9]{2-9} & Производная на отрезке & Максимум функции спутан с поведением производной & $2$ & $-1$ \\
\hyperref[task:2_12]{2-12} & График логарифма & Решение по двум точкам не доведено до ответа & Ответ не указан & $\frac18$ \\
\hyperref[task:2_13]{2-13} & Кредит и сумма платежей & Неправильно посчитан суммарный платёж & $2{,}125$ млн руб. & $1{,}9896$ млн руб. \\
\hyperref[task:2_17]{2-17} & Максимум распределения Максвелла & Решение производной не доведено до экстремума & Ответ не указан & Около $510$ м/с \\
\hyperref[task:3_1]{3-1} & Трапеция и средняя линия & Неверно найдена половина основания & $11$ & $6{,}5$ \\
\hyperref[task:3_4]{3-4} & Классическая вероятность & Неправильное отношение числа исходов & $0{,}45$ & $0{,}2$ \\
\hyperref[task:3_5]{3-5} & Вероятности событий & Перепутаны вероятности противоположных событий & $0{,}6$ & $0{,}14$ \\
\hyperref[task:3_9]{3-9} & Производная и скорость & Найдено расстояние вместо скорости & $14{,}7$ & $3{,}4$ м/с \\
\hyperref[task:3_10]{3-10} & Адиабатическое сжатие & Перевернуто отношение давлений & $76{,}8$ л & $0{,}075$ л \\
\hyperref[task:3_14]{3-14} & Показательное уравнение & Потеряна пара корней и не выполнен отбор & $1$ и $2$ & $1$ на $[0{,}4;1{,}4]$ \\
\hyperref[task:3_16]{3-16} & Неравенство с логарифмами & Неверно преобразованы основания логарифмов; решение оборвано & Ответ не указан & $|x|\le 2\sqrt5$ либо\newline $\frac{\sqrt{95}}2\le|x|$\newline $<2\sqrt6$ \\
\bottomrule
\end{longtable}
\normalsize
Номера вида «1--7» означают: лист 1, задание 7. Пустые пункты, на которых нет следов решения, здесь не оцениваются.
\newpage
\thispagestyle{empty}
\vspace*{\fill}\begin{center}{\LARGE\bfseries Разбор ошибок}\end{center}\vspace*{\fill}
\newpage
\phantomsection\label{task:1_7}
\section*{Лист 1, задание № 7. Логарифмы и ОДЗ}
\Heading{Текст задания}{Решите уравнение $\log_5(0{,}25-x)=\log_5(x^2-x)$. Если корней несколько, укажите больший.}
\Heading{Ваше решение}{Вы приравняли подлогарифмические выражения, получили $x^2=0{,}25$ и два кандидата $x=0{,}5$ и $x=-0{,}5$. В ответ перенесено число $0{,}5$.}
\Heading{Ваш ответ}{$0{,}5$}
\Heading{Ошибка}{\textcolor{red}{\textbf{Ошибка:}} Не учтена область допустимых значений логарифмов.}
\Heading{Где именно возникает ошибка}{Приравнивание аргументов логарифмов с одинаковым основанием --- правильный метод, но он применяется после проверки положительности аргументов.}
\Heading{Первый неверный или недостающий шаг}{Первый неверный шаг --- объявление $0{,}5$ допустимым корнем без подстановки в оба аргумента.}
\Heading{Почему это неверно}{Логарифм действительного числа определён только для положительного аргумента. При $x=0{,}5$ имеем $0{,}25-0{,}5=-0{,}25<0$, и выражение слева не существует. Равенство многочленов ещё не означает равенства логарифмов.}
\Heading{Ключевая идея}{\Idea{} ОДЗ проверяется до окончательного выбора корней; посторонние кандидаты необходимо исключать.}
\Heading{Исправление от последнего правильного шага}{Запишите два условия: $0{,}25-x>0$ и $x^2-x>0$. Совместно они дают $x<0$. Следовательно, кандидат $0{,}5$ исключён.}
\Heading{Полное правильное решение}{Из монотонности логарифма по основанию $5$ получаем $0{,}25-x=x^2-x$, поэтому $x^2=\frac14$, то есть $x=\pm\frac12$. При $x=\frac12$ оба аргумента отрицательны, а при $x=-\frac12$ оба равны $\frac34>0$. Значит, допустим только второй корень.}
\Heading{Проверка}{Подстановка $x=-\frac12$ даёт $\log_5\frac34=\log_5\frac34$, что подтверждает решение.}
\Heading{Итог}{\Res{} $x=-\frac12$.}
\Heading{Задача на закрепление}{Решите уравнение $\log_3(1-x)=\log_3(x^2+x)$.}
\newpage
\phantomsection\label{task:1_9}
\section*{Лист 1, задание № 9. Производная и экстремумы}
\Heading{Текст задания}{По графику $y=f\prime(x)$ определите количество точек максимума функции $f(x)$ на отрезке $[-9;3]$.}
\Heading{Ваше решение}{Вы записали признаки экстремума через производную, отметили границы отрезка и указали в ответе три точки максимума.}
\Heading{Ваш ответ}{$3$}
\Heading{Ошибка}{\textcolor{red}{\textbf{Ошибка:}} Перепутаны нули производной и точки максимума.}
\Heading{Где именно возникает ошибка}{Верно, что нужно анализировать график производной, а не высоты вершин её графика.}
\Heading{Первый неверный или недостающий шаг}{Ошибка возникает при подсчёте: в качестве максимума учтён не только переход производной от положительных значений к отрицательным.}
\Heading{Почему это неверно}{Точка максимума функции $f$ появляется, если при движении слева направо производная $f\prime$ меняет знак с «плюса» на «минус». Обратный переход означает минимум; касание оси без изменения знака вообще не даёт основания называть точку экстремумом.}
\Heading{Ключевая идея}{\Idea{} Каждое пересечение графиком $f\prime$ горизонтальной оси проверяйте по знакам слева и справа.}
\Heading{Исправление от последнего правильного шага}{На данном графике внутри $[-9;3]$ два перехода $f\prime:+\to-$. Остальные отмеченные пересечения либо дают минимум, либо не принадлежат требуемому типу.}
\Heading{Полное правильное решение}{Слева от первого нужного пересечения $f\prime>0$, справа $f\prime<0$: функция сначала возрастает, затем убывает. Аналогичный переход происходит ещё один раз ближе к правой части отрезка. Между ними встречаются переходы $-\to+$, при которых $f$ переходит от убывания к возрастанию; их считать нельзя. Получаем ровно две точки максимума.}
\Heading{Проверка}{Последовательно выписываем знаки на промежутках между нулями: знак «плюс» перед каждым из двух выбранных нулей и «минус» после каждого согласуется с видом графика.}
\Heading{Итог}{\Res{} Количество точек максимума равно двум.}
\Heading{Задача на закрепление}{Для функции $g$ на $[-4;4]$ знаки $g\prime$ на последовательных промежутках между нулями $-2$, $0$, $3$ равны $+$, $-$, $+$, $-$. Найдите число точек максимума.}
\newpage
\phantomsection\label{task:1_12}
\section*{Лист 1, задание № 12. Корни и графики}
\Heading{Текст задания}{На рисунке даны графики $f(x)=\sqrt{ax}$ и $g(x)=\sqrt{x-b}$. Найдите абсциссу их точки пересечения.}
\Heading{Ваше решение}{На рисунке Вы отметили несколько точек и выписали подстановки координат в формулы корней; итоговое значение абсциссы не получено.}
\Heading{Ваш ответ}{Ответ не указан}
\Heading{Ошибка}{\textcolor{red}{\textbf{Ошибка:}} Не установлен параметр каждого графика и не доведено до конца уравнение пересечения.}
\Heading{Где именно возникает ошибка}{Верно, что координаты точек помогают восстановить параметры функций.}
\Heading{Первый неверный или недостающий шаг}{Далее координаты разных кривых нельзя подставлять без проверки: кривая $\sqrt{ax}$ проходит через начало координат, а $\sqrt{x-b}$ начинается в точке $(b;0)$.}
\Heading{Почему это неверно}{Изображены две различные корневые функции. Первая, начинающаяся при $x=0$, задаёт параметр $a$; вторая, начинающаяся при отрицательном $x$, задаёт параметр $b$. Смешивание этих кривых приводит к неверным параметрам и точке пересечения.}
\Heading{Ключевая идея}{\Idea{} Сначала определите, какой формуле принадлежит каждая кривая, затем используйте точки с точными координатами.}
\Heading{Исправление от последнего правильного шага}{Для правой кривой точка $(3;2)$ даёт $2=\sqrt{3a}$, поэтому $a=4/3$. Для левой кривой точка $(-1;2)$ даёт $2=\sqrt{-1-b}$, поэтому $b=-5$.}
\Heading{Полное правильное решение}{Ищем пересечение: $\sqrt{\frac43x}=\sqrt{x+5}$. Обе стороны неотрицательны при $x\ge0$, значит, допустимо возвести их в квадрат: $\frac43x=x+5$, откуда $x/3=5$ и $x=15$.}
\Heading{Проверка}{При $x=15$ получаем $\sqrt{\frac43\cdot15}=\sqrt{20}$ и $\sqrt{15+5}=\sqrt{20}$; равенство выполнено, ОДЗ соблюдена.}
\Heading{Итог}{\Res{} Абсцисса пересечения равна $15$.}
\Heading{Задача на закрепление}{Функции $y=\sqrt{2x}$ и $y=\sqrt{x+6}$ пересекаются в одной точке. Найдите её абсциссу.}
\newpage
\phantomsection\label{task:1_13}
\section*{Лист 1, задание № 13. Кредит по годам}
\Heading{Текст задания}{В июле 2030 года выдан кредит $1{,}5$ млн рублей на четыре года под $10\%$ годовых. В январе долг возрастает на $10\%$, с февраля по июнь идут выплаты, в июле 2031--2033 годов долг на $40\%$ меньше, чем в июле предыдущего года; в июле 2034 года долг погашен. Найдите платёж 2034 года.}
\Heading{Ваше решение}{В ответе указано $1{,}8264$ млн рублей --- сумма платежей за все четыре года, а не платёж последнего года.}
\Heading{Ваш ответ}{$1{,}8264$ млн руб.}
\Heading{Ошибка}{\textcolor{red}{\textbf{Ошибка:}} Подмена требуемой величины.}
\Heading{Где именно возникает ошибка}{Можно правильно рассчитать график погашения при помощи июльских остатков долга.}
\Heading{Первый неверный или недостающий шаг}{Ошибка в заключении: сумма четырёх ежегодных платежей названа платежом 2034 года.}
\Heading{Почему это неверно}{По условию требуется платёж лишь в последний год. К январю 2034-го процентов начисляется только на долг, оставшийся в июле 2033-го. Уже внесённые платежи 2031--2033 годов не нужно складывать с последним платежом.}
\Heading{Ключевая идея}{\Idea{} Разделяйте «остаток долга», «платёж за конкретный год» и «все платежи вместе».}
\Heading{Исправление от последнего правильного шага}{Июльские остатки после платежей: $1{,}5$; $0{,}9$; $0{,}54$; $0{,}324$ млн рублей. Последний остаток увеличивается в январе на $10\%$.}
\Heading{Полное правильное решение}{Пусть $D_k$ --- долг в июле года. Тогда $D_{2031}=0{,}6\cdot1{,}5=0{,}9$, $D_{2032}=0{,}6\cdot0{,}9=0{,}54$, $D_{2033}=0{,}6\cdot0{,}54=0{,}324$ млн. В январе 2034 года долг станет $1{,}1\cdot0{,}324=0{,}3564$ млн. До июля он должен быть полностью выплачен, следовательно, платёж равен $356400$ рублям.}
\Heading{Проверка}{Можно отдельно проверить всю последовательность: выплаты $0{,}75$, $0{,}45$, $0{,}27$ и $0{,}3564$ млн рублей. Их сумма $1{,}8264$ млн совпадает с записанным числом, но это другая величина.}
\Heading{Итог}{\Res{} В 2034 году выплачено $356400$ рублей.}
\Heading{Задача на закрепление}{Кредит $2$ млн руб. взят в июле на четыре года под $12\%$ годовых. После платежей в июле каждого из первых трёх лет остаток долга сокращается на $50\%$ относительно прошлого июля. Определите платёж четвёртого года.}
\newpage
\phantomsection\label{task:1_14}
\section*{Лист 1, задание № 14. Тригонометрическое уравнение}
\Heading{Текст задания}{Решите $24\sin2x-7\cos2x=15$ и найдите корни на $[-3\pi;-2\pi]$.}
\Heading{Ваше решение}{После верного раскрытия формул двойного угла Вы получили выражение $48\sin x\cos x-22\cos^2x-8\sin^2x=0$, но затем представили его произведением, не равным исходной сумме. Из этого получены неверные серии корней.}
\Heading{Ваш ответ}{Получены посторонние серии корней}
\Heading{Ошибка}{\textcolor{red}{\textbf{Ошибка:}} Неверное алгебраическое разложение на множители.}
\Heading{Где именно возникает ошибка}{Шаг с приведением $15$ к $15(\sin^2x+\cos^2x)$ и раскрытием формул двойного угла выполнен правильно.}
\Heading{Первый неверный или недостающий шаг}{Первый ошибочный переход --- замена $48\sin x\cos x-22\cos^2x-8\sin^2x$ на произведение, где теряются квадратичные слагаемые.}
\Heading{Почему это неверно}{Раскрытие предложенных скобок не восстанавливает исходный трёхчлен. Нельзя сокращать или выносить множитель, который не является общим для всех трёх членов. Разложение нужно обязательно проверять обратным раскрытием скобок.}
\Heading{Ключевая идея}{\Idea{} Выражение однородно по $\sin x$ и $\cos x$. Проверив, что $\cos x\ne0$, можно разделить на $\cos^2x$ и перейти к $t=\tan x$.}
\Heading{Исправление от последнего правильного шага}{При $\cos x=0$ исходная левая часть равна $7$, а не $15$, поэтому такие значения не подходят. Делим правильное равенство на $\cos^2x$: $48\tan x-22-8\tan^2x=0$.}
\Heading{Полное правильное решение}{Пусть $t=\tan x$. Получаем $4t^2-24t+11=0$. Дискриминант $D=576-176=400$, поэтому $t=(24\pm20)/8$. Значит, $t_1=1/2$, $t_2=11/2$. Общее решение: $x=\arctan(1/2)+\pi n$ или $x=\arctan(11/2)+\pi n$. Оба арктангенса лежат между $0$ и $\pi/2$; в отрезке $[-3\pi;-2\pi]$ для каждой серии подходит только $n=-3$.}
\Heading{Проверка}{Оба найденных значения удовлетворяют уравнению $4t^2-24t+11=0$, эквивалентному исходному при $\cos x\ne0$. По расположению на числовой оси оба лежат на заданном отрезке.}
\Heading{Итог}{\Res{} $x=-3\pi+\arctan\frac12$ и $x=-3\pi+\arctan\frac{11}{2}$.}
\Heading{Задача на закрепление}{Решите $10\sin2x-3\cos2x=5$ и выпишите корни из промежутка $[-\pi;0]$.}
\newpage
\phantomsection\label{task:2_1}
\section*{Лист 2, задание № 1. Площадь треугольника}
\Heading{Текст задания}{Угол равнобедренного треугольника при вершине равен $30^\circ$, боковая сторона равна $\sqrt6$. Найдите площадь.}
\Heading{Ваше решение}{Вы правильно обозначили равные боковые стороны и нашли углы при основании по $75^\circ$, однако записали площадь $\frac12$.}
\Heading{Ваш ответ}{$\frac12$}
\Heading{Ошибка}{\textcolor{red}{\textbf{Ошибка:}} Не применена формула площади через две стороны и угол между ними.}
\Heading{Где именно возникает ошибка}{Углы при основании действительно равны $(180^\circ-30^\circ)/2=75^\circ$.}
\Heading{Первый неверный или недостающий шаг}{Первый недостающий шаг --- вычисление площади по известным двум сторонам длиной $\sqrt6$ и заключённому между ними углу $30^\circ$.}
\Heading{Почему это неверно}{Площадь не равна просто половине произведения неопределённых длин. В формуле обязательно учитывается синус угла между сторонами. Для угла $30^\circ$ он равен $1/2$, но произведение боковых сторон равно $6$, а не $2$.}
\Heading{Ключевая идея}{\Idea{} Если известны две стороны и угол между ними, сразу применяйте $S=\frac12 ab\sin\gamma$.}
\Heading{Исправление от последнего правильного шага}{Берём $a=b=\sqrt6$, $\gamma=30^\circ$, не проводя дополнительной высоты.}
\Heading{Полное правильное решение}{$S=\frac12\sqrt6\cdot\sqrt6\cdot\sin30^\circ=\frac12\cdot6\cdot\frac12=\frac32$.}
\Heading{Проверка}{Результат меньше площади прямоугольника со сторонами $\sqrt6$, то есть меньше $6$, что согласуется с рисунком. Проверка через высоту даёт ту же величину.}
\Heading{Итог}{\Res{} Площадь равна $\frac32$.}
\Heading{Задача на закрепление}{Угол между равными сторонами равнобедренного треугольника равен $60^\circ$, каждая боковая сторона имеет длину $4$. Найдите площадь.}
\newpage
\phantomsection\label{task:2_5}
\section*{Лист 2, задание № 5. Вероятность двух матчей}
\Heading{Текст задания}{Команда А играет дома и в гостях. Дома победа имеет вероятность $0{,}6$, ничья $0{,}1$; в гостях победа $0{,}3$, ничья $0{,}2$. За победу дают три очка, за ничью одно. Найдите вероятность набрать не менее четырёх очков.}
\Heading{Ваше решение}{Вы выписали отдельные вероятности исходов и два числа без однозначного итогового расчёта; суммы вероятностей необходимых совместных событий нет.}
\Heading{Ваш ответ}{На листе записаны $0{,}3$ и $0{,}5$}
\Heading{Ошибка}{\textcolor{red}{\textbf{Ошибка:}} Не завершено разбиение события «не менее четырёх очков» на случаи.}
\Heading{Где именно возникает ошибка}{Правильно выделено, что в каждом матче три возможных исхода и за них начисляются разные очки.}
\Heading{Первый неверный или недостающий шаг}{Не записаны произведения вероятностей для всех пар результатов, обеспечивающих четыре или шесть очков.}
\Heading{Почему это неверно}{Для четырёх очков нужна победа и ничья в любом порядке, для шести --- две победы. Одна победа и поражение дают только три очка. Поскольку матчи независимы по условию, вероятности отдельных сочетаний перемножают, а несовместимые варианты складывают.}
\Heading{Ключевая идея}{\Idea{} Сначала перечислите подходящие пары очков: $(3,3)$, $(3,1)$ и $(1,3)$.}
\Heading{Исправление от последнего правильного шага}{Вероятности трёх сочетаний равны $0{,}6\cdot0{,}3$, $0{,}6\cdot0{,}2$ и $0{,}1\cdot0{,}3$.}
\Heading{Полное правильное решение}{Вычисляем: $P(3,3)=0{,}18$; $P(3,1)=0{,}12$; $P(1,3)=0{,}03$. Эти события несовместимы, поэтому $P=0{,}18+0{,}12+0{,}03=0{,}33$.}
\Heading{Проверка}{Сумма положительна и меньше единицы; все допустимые пары перечислены, других способов набрать четыре или более очков за два матча нет.}
\Heading{Итог}{\Res{} Искомая вероятность $0{,}33$.}
\Heading{Задача на закрепление}{Дома вероятность победы $0{,}5$, ничьей $0{,}2$; в гостях победы $0{,}4$, ничьей $0{,}1$. Найдите вероятность получить не менее четырёх очков за два независимых матча.}
\newpage
\phantomsection\label{task:2_8}
\section*{Лист 2, задание № 8. Логарифмы и дроби}
\Heading{Текст задания}{Вычислите $\log_a(b^6\sqrt a)$, если $\log_b a=-2$.}
\Heading{Ваше решение}{Вы воспользовались формулой перехода к обратному логарифму, получили $\log_a b=-\frac12$ и дошли до $6\cdot(-\frac12)+\frac12$. При сложении записали $-\frac72$.}
\Heading{Ваш ответ}{$-3{,}5$}
\Heading{Ошибка}{\textcolor{red}{\textbf{Ошибка:}} Неверно сложены числа $-3$ и $\frac12$.}
\Heading{Где именно возникает ошибка}{Преобразования $\log_a b=1/\log_b a=-1/2$ и $\log_a\sqrt a=1/2$ выполнены верно.}
\Heading{Первый неверный или недостающий шаг}{Первый неверный переход: $-3+\frac12=-\frac72$.}
\Heading{Почему это неверно}{Чтобы сложить целое число и половину, представляем $-3=-6/2$. Тогда $-6/2+1/2=-5/2$. Знак второго слагаемого положительный; менять его на отрицательный нельзя.}
\Heading{Ключевая идея}{\Idea{} Отдельно выполняйте преобразования логарифмов, отдельно --- арифметику дробей со знаками.}
\Heading{Исправление от последнего правильного шага}{Возвращаемся к Вашей правильной строке $6(-\frac12)+\frac12$ и аккуратно складываем.}
\Heading{Полное правильное решение}{$\log_a(b^6\sqrt a)=6\log_a b+\frac12\log_a a=6(-\frac12)+\frac12=-3+\frac12=-\frac52=-2{,}5$. Предполагается, что основания логарифмов допустимы, как и в условии.}
\Heading{Проверка}{$-\frac52=-2{,}5$ и $-\frac72=-3{,}5$ различаются ровно на единицу, что подтверждает локальность арифметической ошибки.}
\Heading{Итог}{\Res{} Значение выражения равно $-\frac52$.}
\Heading{Задача на закрепление}{Найдите $\log_a(c^4\sqrt a)$, если $\log_c a=-4$.}
\newpage
\phantomsection\label{task:2_9}
\section*{Лист 2, задание № 9. Производная на отрезке}
\Heading{Текст задания}{По графику $y=f\prime(x)$ найдите абсциссу точки максимума функции $f$ на отрезке $[-2;6]$.}
\Heading{Ваше решение}{В ответ записано $2$; рядом верно указано, что надо обратиться к производной, но не проверен знак $f\prime(2)$.}
\Heading{Ваш ответ}{$2$}
\Heading{Ошибка}{\textcolor{red}{\textbf{Ошибка:}} Неверно определена точка максимума по графику производной.}
\Heading{Где именно возникает ошибка}{Верное направление поиска --- анализ знака $f\prime$, а не высоты её графика.}
\Heading{Первый неверный или недостающий шаг}{При $x=2$ производная отрицательна и не меняет знак с плюса на минус, поэтому эта точка не может быть найденным максимумом.}
\Heading{Почему это неверно}{Если $f\prime(x)<0$ в окрестности точки, функция убывает. Точка максимума на данном графике соответствует переходу производной из положительных значений в отрицательные возле $x=-1$. Справа у функции есть участок возрастания, но он не компенсирует предыдущее уменьшение до высоты максимума.}
\Heading{Ключевая идея}{\Idea{} Проверяйте характер изменения самой функции: $f\prime>0$ означает рост, $f\prime<0$ --- убывание.}
\Heading{Исправление от последнего правильного шага}{На $[-2;-1]$ производная положительна; затем до следующего нуля около $x=5$ отрицательна; после него до $6$ положительна. Кандидаты для наибольшего значения --- $x=-1$ и граничная точка $x=6$.}
\Heading{Полное правильное решение}{По знаку и площади участков под графиком $f\prime$ видно, что накопленное убывание после $-1$ существенно больше заключительного прироста перед $6$. Следовательно, $f(-1)>f(6)$, а на первом участке $f$ возрастает к точке $-1$. Абсцисса точки наибольшего значения $-1$.}
\Heading{Проверка}{При $x=-1$ на графике $f\prime$ меняет знак $+\to-$, что подтверждает наличие максимума.}
\Heading{Итог}{\Res{} Искомая абсцисса равна $-1$.}
\Heading{Задача на закрепление}{Пусть $f\prime(x)=-(x+1)(x-3)$ на $[-2;4]$. Найдите абсциссу точки максимума функции $f$.}
\newpage
\phantomsection\label{task:2_12}
\section*{Лист 2, задание № 12. Логарифмический график}
\Heading{Текст задания}{По графику функции $f(x)=k\log_a x$ найдите $x$, если $f(x)=-9$. На графике отмечены точки $(2;3)$ и $(4;6)$.}
\Heading{Ваше решение}{Вы переписали координаты двух точек и составили равенства для неизвестных параметров, но значение $x$ не получено.}
\Heading{Ваш ответ}{Ответ не указан}
\Heading{Ошибка}{\textcolor{red}{\textbf{Ошибка:}} Не использовано свойство логарифма для изменения аргумента в степени.}
\Heading{Где именно возникает ошибка}{Верно замечены значения $f(2)=3$ и $f(4)=6$. Второе согласуется с тем, что $4=2^2$.}
\Heading{Первый неверный или недостающий шаг}{Решение оборвано до определения аргумента, при котором значение логарифмической функции становится равным $-9$.}
\Heading{Почему это неверно}{Логарифм степени равен показателю, умноженному на логарифм основания этой степени. Поэтому $f(2^m)=m f(2)$ независимо от того, какие именно $k$ и $a$ определяют график (при допустимых параметрах).}
\Heading{Ключевая идея}{\Idea{} Удобно сравнивать значения при аргументах $2$, $2^2$, $2^{-3}$ вместо поиска обоих параметров по отдельности.}
\Heading{Исправление от последнего правильного шага}{Так как $f(2)=3$, для $f(x)=-9=-3\cdot3$ достаточно взять $x=2^{-3}$.}
\Heading{Полное правильное решение}{$f(2^{-3})=k\log_a(2^{-3})=-3k\log_a2=-3f(2)=-9$. Следовательно, $x=2^{-3}=\frac18$.}
\Heading{Проверка}{Аргумент $1/8$ положителен, значит, логарифм определён; $f(1/8)=-9$ доказано подстановкой.}
\Heading{Итог}{\Res{} Искомый аргумент $x=\frac18$.}
\Heading{Задача на закрепление}{Функция $y=c\log_b x$ проходит через точку $(3;4)$. Найдите $x$, при котором $y=-8$.}
\newpage
\phantomsection\label{task:2_13}
\section*{Лист 2, задание № 13. Кредит и сумма платежей}
\Heading{Текст задания}{В июле 2030 года взят кредит $1{,}5$ млн рублей на четыре года под $15\%$ годовых. В январе начисляются проценты; к июлю 2031--2033 годов долг уменьшают на $40\%$ относительно предыдущего июля; к июлю 2034 года кредит погашают полностью. Найдите сумму всех платежей.}
\Heading{Ваше решение}{В ответ записано $2{,}125$ млн рублей; согласованный с календарём график остатков и выплат не приведён.}
\Heading{Ваш ответ}{$2{,}125$ млн руб.}
\Heading{Ошибка}{\textcolor{red}{\textbf{Ошибка:}} Сумма выплат не соответствует заданному графику погашения.}
\Heading{Где именно возникает ошибка}{Правильно выбран смысл величины: здесь в отличие от первого листа нужна сумма платежей за все годы.}
\Heading{Первый неверный или недостающий шаг}{Не учтено, что после ежегодного сокращения остатка долга проценты следующего года начисляются на меньшую сумму.}
\Heading{Почему это неверно}{Проценты начисляются в январе на остаток долга предыдущего июля. Платёж каждого из первых трёх лет равен долгу после начисления процентов минус запланированный новый июльский остаток; последний платёж равен всему долгу января 2034-го.}
\Heading{Ключевая идея}{\Idea{} Сначала сформируйте последовательность июльских остатков, затем вычисляйте каждый годовой платёж.}
\Heading{Исправление от последнего правильного шага}{Остатки долга в июле 2030, 2031, 2032, 2033 годов: $1{,}5$, $0{,}9$, $0{,}54$, $0{,}324$ млн. В январе они умножаются на $1{,}15$.}
\Heading{Полное правильное решение}{Платёж 2031 года: $1{,}15\cdot1{,}5-0{,}9=0{,}825$ млн. Платёж 2032-го: $1{,}15\cdot0{,}9-0{,}54=0{,}495$ млн. Платёж 2033-го: $1{,}15\cdot0{,}54-0{,}324=0{,}297$ млн. В 2034-м платёж $1{,}15\cdot0{,}324=0{,}3726$ млн. Сумма $0{,}825+0{,}495+0{,}297+0{,}3726=1{,}9896$ млн рублей.}
\Heading{Проверка}{Общий платёж больше суммы кредита на $0{,}4896$ млн --- это суммарные проценты. Проверяем отдельно: $0{,}15(1{,}5+0{,}9+0{,}54+0{,}324)=0{,}4896$.}
\Heading{Итог}{\Res{} Общая сумма платежей $1\,989\,600$ рублей.}
\Heading{Задача на закрепление}{На четыре года взят кредит $1$ млн рублей под $10\%$ годовых. В июле первых трёх лет остаток уменьшается вдвое относительно предыдущего июля, в четвёртый год погашается полностью. Найдите сумму платежей.}
\newpage
\phantomsection\label{task:2_17}
\section*{Лист 2, задание № 17. Максимум распределения Максвелла}
\Heading{Текст задания}{Плотность скоростей молекул кислорода задана формулой Максвелла\[f(v)=4\pi v^2\left(\frac{M}{2\pi RT}\right)^{3/2}e^{-Mv^2/(2RT)}.\]При $T=500\ \text{К}$ и $M=0{,}032\ \text{кг/моль}$ найдите скорость максимума.}
\Heading{Ваше решение}{В тетради произведена подстановка констант и начато дифференцирование экспоненциального множителя, но не составлено уравнение $f\prime(v)=0$ и не найдено числовое значение.}
\Heading{Ваш ответ}{Ответ не указан}
\Heading{Ошибка}{\textcolor{red}{\textbf{Ошибка:}} Не завершена проверка максимума произведения степенной и показательной функций.}
\Heading{Где именно возникает ошибка}{Вы правильно выделили показатель экспоненты $-Mv^2/(2RT)$ и начали искать производную.}
\Heading{Первый неверный или недостающий шаг}{Решение останавливается до применения правила производной произведения $v^2 e^{-cv^2}$.}
\Heading{Почему это неверно}{Здесь нельзя искать максимум только одного множителя: $v^2$ возрастает, а экспонента при $v>0$ убывает. Максимум их произведения находится через производную всего выражения. Числовой положительный множитель не влияет на положение максимума.}
\Heading{Ключевая идея}{\Idea{} Обозначьте положительные константы одной буквой, найдите производную и исследуйте её знак.}
\Heading{Исправление от последнего правильного шага}{Пусть $C=4\pi\left(\frac{M}{2\pi RT}\right)^{3/2}>0$, $c=M/(2RT)>0$. Тогда $f(v)=Cv^2e^{-cv^2}$.}
\Heading{Полное правильное решение}{$f\prime(v)=C(2ve^{-cv^2}-2cv^3e^{-cv^2})=2Cve^{-cv^2}(1-cv^2)$. При $v>0$ первые множители положительны, поэтому знак зависит от $1-cv^2$. Он меняется с плюса на минус при $v^2=1/c=2RT/M$. Следовательно, $v_*=\sqrt{2RT/M}$. Подставляем $R\approx8{,}31$, $T=500$, $M=0{,}032$: $v_*\approx\sqrt{259687{,}5}\approx510$ м/с.}
\Heading{Проверка}{Слева от найденной положительной скорости производная положительна, справа отрицательна; это действительно максимум. При $v=0$ плотность равна нулю.}
\Heading{Итог}{\Res{} Скорость максимума примерно $510$ м/с.}
\Heading{Задача на закрепление}{Для плотности $g(v)=Av^2e^{-v^2/200}$ при $v\ge0$ и $A>0$ найдите скорость, при которой $g$ максимальна.}
\newpage
\phantomsection\label{task:3_1}
\section*{Лист 3, задание № 1. Трапеция и средняя линия}
\Heading{Текст задания}{Основания трапеции равны $13$ и $19$. Найдите меньший из отрезков, на которые диагональ делит среднюю линию трапеции.}
\Heading{Ваше решение}{Вы записали $13+19=22$, а затем $22/2=11$. Оба основания определены верно, но первое сложение ошибочно.}
\Heading{Ваш ответ}{$11$}
\Heading{Ошибка}{\textcolor{red}{\textbf{Ошибка:}} Не использовано свойство деления средней линии диагональю.}
\Heading{Где именно возникает ошибка}{Верно, что длины средней линии и её частей определяются основаниями трапеции.}
\Heading{Первый неверный или недостающий шаг}{Первый неверный вычислительный шаг --- $13+19=22$, тогда как $13+19=32$. Кроме того, даже правильная полусумма даёт длину всей средней линии, а не её меньшей части.}
\Heading{Почему это неверно}{Диагональ образует два треугольника, отрезки средней линии которых параллельны основаниям трапеции. Каждый из этих отрезков составляет половину соответствующего основания. Поэтому длины частей --- $13/2$ и $19/2$, а не половина суммы двух оснований.}
\Heading{Ключевая идея}{\Idea{} Диагональ разбивает среднюю линию на половины двух оснований.}
\Heading{Исправление от последнего правильного шага}{Меньшее основание равно $13$, значит меньшая часть средней линии $13/2$.}
\Heading{Полное правильное решение}{Обозначим основания $a=13$ и $b=19$. По теореме о средней линии треугольника диагональ пересекает среднюю линию трапеции так, что образуются части длиной $a/2$ и $b/2$. Следовательно, $m_1=13/2=6{,}5$, $m_2=19/2=9{,}5$. Меньшая часть --- $6{,}5$.}
\Heading{Проверка}{Их сумма $6{,}5+9{,}5=16$ совпадает со средней линией трапеции $(13+19)/2=16$.}
\Heading{Итог}{\Res{} Меньший отрезок равен $6{,}5$.}
\Heading{Задача на закрепление}{Основания трапеции равны $10$ и $24$. На какие отрезки диагональ делит её среднюю линию? Найдите меньший.}
\newpage
\phantomsection\label{task:3_4}
\section*{Лист 3, задание № 4. Классическая вероятность}
\Heading{Текст задания}{Среди натуральных чисел от $12$ до $36$ включительно случайно выбирают одно. Найдите вероятность того, что выбранное число делится на $6$.}
\Heading{Ваше решение}{В решении перечислены подходящие числа и составлена дробь, однако записан результат $0{,}45$.}
\Heading{Ваш ответ}{$0{,}45$}
\Heading{Ошибка}{\textcolor{red}{\textbf{Ошибка:}} Неверно подсчитано общее число равновозможных исходов или вычислена дробь.}
\Heading{Где именно возникает ошибка}{Правильно выделено событие «число кратно шести».}
\Heading{Первый неверный или недостающий шаг}{Ответ $0{,}45$ не получается из точного количества кратных и общего числа целых чисел на указанном отрезке.}
\Heading{Почему это неверно}{Числа $12$ и $36$ включены, поэтому всего вариантов $36-12+1=25$. Кратные шести: $12$, $18$, $24$, $30$, $36$ --- их пять. Вероятность определяется отношением благоприятных исходов ко всем равновозможным.}
\Heading{Ключевая идея}{\Idea{} При счёте целых чисел на отрезке включайте обе границы: количество равно «правый минус левый плюс один».}
\Heading{Исправление от последнего правильного шага}{Подставляем $m=5$ и $n=25$ в формулу $P=m/n$.}
\Heading{Полное правильное решение}{$P=\frac5{25}=\frac15=0{,}2$.}
\Heading{Проверка}{Проверка: приблизительно каждое пятое целое число в таком наборе кратно шести; результат $20\%$ согласуется с перечнем.}
\Heading{Итог}{\Res{} Вероятность равна $0{,}2$.}
\Heading{Задача на закрепление}{Из целых чисел от $8$ до $32$ включительно выбирают одно. Найдите вероятность того, что оно кратно $4$.}
\newpage
\phantomsection\label{task:3_5}
\section*{Лист 3, задание № 5. Вероятности событий}
\Heading{Текст задания}{Вероятность температуры выше $+22^\circ\text{C}$ равна $0{,}27$, а вероятность температуры ниже $+25^\circ\text{C}$ равна $0{,}87$. Найдите вероятность температуры в промежутке от $+22^\circ\text{C}$ до $+25^\circ\text{C}$.}
\Heading{Ваше решение}{Вы вычислили разность $0{,}87-0{,}27=0{,}6$.}
\Heading{Ваш ответ}{$0{,}6$}
\Heading{Ошибка}{\textcolor{red}{\textbf{Ошибка:}} Из вероятности верхней границы вычтено не то событие.}
\Heading{Где именно возникает ошибка}{Правильная идея --- выразить вероятность попадания в интервал через разность двух накопленных вероятностей.}
\Heading{Первый неверный или недостающий шаг}{Первый неверный шаг: из $P(T<25)$ вычтено $P(T>22)$ вместо вероятности события $T\le22$.}
\Heading{Почему это неверно}{Чтобы оставить температуры между $22$ и $25$, нужно из события $T<25$ убрать все температуры, не превышающие $22$. Это противоположное событие к $T>22$, поэтому его вероятность равна $1-0{,}27=0{,}73$. Простая разность $0{,}87-0{,}27$ сравнивает события, которые не вложены требуемым образом.}
\Heading{Ключевая идея}{\Idea{} Сначала изобразите нужный интервал и укажите, какую часть большего события следует убрать.}
\Heading{Исправление от последнего правильного шага}{$P(T\le22)=1-P(T>22)=0{,}73$. Затем $P(22<T<25)=0{,}87-0{,}73=0{,}14$ (для измерения непрерывной температуры вероятность попадания точно на границу равна нулю).}
\Heading{Полное правильное решение}{Используем свойства вероятности: $P(T<25)=P(T\le22)+P(22<T<25)$. Подставив известные числа, имеем $0{,}87=0{,}73+p$. Значит, $p=0{,}14$.}
\Heading{Проверка}{Подставив найденное число обратно, получим $0{,}73+0{,}14=0{,}87$. Значение находится между нулём и единицей.}
\Heading{Итог}{\Res{} Вероятность равна $0{,}14$.}
\Heading{Задача на закрепление}{Вероятность, что температура выше $10^\circ\text{C}$, равна $0{,}6$, а ниже $18^\circ\text{C}$ --- $0{,}85$. Найдите вероятность, что она между этими значениями.}
\newpage
\phantomsection\label{task:3_9}
\section*{Лист 3, задание № 9. Производная и скорость}
\Heading{Текст задания}{Материальная точка движется по закону $x(t)=0{,}2t^2+t+1{,}5$, где $x$ --- метры, $t$ --- секунды. Найдите скорость в момент $t=6$ секунд.}
\Heading{Ваше решение}{Вы подставили $t=6$ прямо в функцию координаты: $x(6)=0{,}2\cdot36+6+1{,}5=14{,}7$.}
\Heading{Ваш ответ}{$14{,}7$}
\Heading{Ошибка}{\textcolor{red}{\textbf{Ошибка:}} Вместо скорости вычислена координата точки.}
\Heading{Где именно возникает ошибка}{Подстановка в формулу координаты выполнена арифметически верно.}
\Heading{Первый неверный или недостающий шаг}{Первый ошибочный шаг --- отождествление $x(6)$ со скоростью в момент $t=6$.}
\Heading{Почему это неверно}{Координата измеряется в метрах, скорость --- в метрах в секунду. Мгновенная скорость прямолинейного движения равна производной координаты по времени: $v(t)=x\prime(t)$. Подстановка времени в $x(t)$ отвечает на другой вопрос: где находится точка в данный момент.}
\Heading{Ключевая идея}{\Idea{} Сначала найдите производную $x\prime(t)$, потом подставьте требуемое время.}
\Heading{Исправление от последнего правильного шага}{$x\prime(t)=(0{,}2t^2)\prime+(t)\prime+(1{,}5)\prime=0{,}4t+1$.}
\Heading{Полное правильное решение}{При $t=6$ получаем $v(6)=0{,}4\cdot6+1=2{,}4+1=3{,}4$ м/с.}
\Heading{Проверка}{По размерностям всё согласуется: производная метров по секундам измеряется в метрах в секунду. Значение координаты $14{,}7$ м не подменяет скорость.}
\Heading{Итог}{\Res{} Скорость равна $3{,}4$ м/с.}
\Heading{Задача на закрепление}{Точка движется по закону $x(t)=0{,}5t^2+3t-2$. Найдите её мгновенную скорость в момент $t=4$ секунды.}
\newpage
\phantomsection\label{task:3_10}
\section*{Лист 3, задание № 10. Адиабатическое сжатие}
\Heading{Текст задания}{Для газа справедливо $p_1V_1^{1{,}4}=p_2V_2^{1{,}4}$. При $V_1=2{,}4$ л, $p_1=0{,}5$ атм найдите объём при $p_2=64$ атм.}
\Heading{Ваше решение}{Вы правильно выписали уравнение адиабаты, но при переносе давлений получили $V_2=32\cdot2{,}4=76{,}8$ л.}
\Heading{Ваш ответ}{$76{,}8$ л}
\Heading{Ошибка}{\textcolor{red}{\textbf{Ошибка:}} Отношение давлений перевернуто при выражении конечного объёма.}
\Heading{Где именно возникает ошибка}{Начальная формула $p_1V_1^{1{,}4}=p_2V_2^{1{,}4}$ выбрана верно.}
\Heading{Первый неверный или недостающий шаг}{Первый неверный переход происходит при изоляции $V_2^{1{,}4}$: отношение $64/0{,}5=128$ должно оказаться в знаменателе, а не в числителе.}
\Heading{Почему это неверно}{Когда давление при сжатии растёт, объём должен уменьшаться. Из формулы следует $V_2^{1{,}4}=(p_1/p_2)V_1^{1{,}4}$. Использование обратного отношения описывало бы увеличение объёма и противоречило самому слову «сжатие».}
\Heading{Ключевая идея}{\Idea{} Сначала выразите искомую степень с символическими давлениями, затем подставляйте числа и контролируйте физический смысл.}
\Heading{Исправление от последнего правильного шага}{$V_2^{1{,}4}=\frac{0{,}5}{64}(2{,}4)^{1{,}4}=\frac1{128}(2{,}4)^{1{,}4}$.}
\Heading{Полное правильное решение}{Отсюда $V_2=2{,}4\left(\frac1{128}\right)^{1/1{,}4}$. Поскольку $1{,}4=7/5$, а $128=2^7$, имеем $V_2=2{,}4(2^{-7})^{5/7}=2{,}4\cdot2^{-5}=\frac{2{,}4}{32}=0{,}075$ л.}
\Heading{Проверка}{Подстановка в исходную формулу: $64(0{,}075)^{1{,}4}=0{,}5(2{,}4)^{1{,}4}$. Также $0{,}075<2{,}4$, что соответствует повышению давления при сжатии.}
\Heading{Итог}{\Res{} Конечный объём $0{,}075$ л.}
\Heading{Задача на закрепление}{Газ сжимают по закону $pV^{1{,}5}=\mathrm{const}$. Давление выросло в $8$ раз, а начальный объём был $4$ л. Найдите конечный объём.}
\newpage
\phantomsection\label{task:3_14}
\section*{Лист 3, задание № 14. Показательное уравнение}
\Heading{Текст задания}{Решите уравнение $4^{x^3-x^2+3x-0{,}5}+4^{x^3+x^2-3x+2{,}5}-5\cdot2^{2x^3+1}=0$ и укажите корни из $[0{,}4;1{,}4]$.}
\Heading{Ваше решение}{Вы свели выражение к квадратному уравнению относительно $t=4^{x^2-3x}>0$, нашли $t=\frac14$ и $t=\frac1{16}$. Затем оба случая фактически заменены уравнением с корнями $1$ и $2$.}
\Heading{Ваш ответ}{$1$ и $2$}
\Heading{Ошибка}{\textcolor{red}{\textbf{Ошибка:}} При возврате от замены к $x$ неверно решено одно из квадратных уравнений; дополнительно не произведён отбор корней по отрезку.}
\Heading{Где именно возникает ошибка}{Сама замена $t=4^{x^2-3x}$, условие $t>0$ и оба значения $t$ найдены правильно.}
\Heading{Первый неверный или недостающий шаг}{При $t=1/4$ нужно решить $x^2-3x=-1$, а не $x^2-3x=-2$. Именно здесь потеряны корни с радикалом.}
\Heading{Почему это неверно}{Две допустимые величины $t$ порождают два разных квадратных уравнения. У первого дискриминант равен $9-4=5$, у второго $9-8=1$. Их нельзя считать одинаковыми; после получения всех корней необходимо проверить принадлежность интервалу $[0{,}4;1{,}4]$.}
\Heading{Ключевая идея}{\Idea{} Возврат из замены --- самостоятельный этап: для каждого значения новой переменной нужно написать отдельное уравнение и выполнить ограничение из пункта «б».}
\Heading{Исправление от последнего правильного шага}{При $t=1/4=4^{-1}$ имеем $x^2-3x=-1$. При $t=1/16=4^{-2}$ имеем $x^2-3x=-2$.}
\Heading{Полное правильное решение}{Делим исходное уравнение на положительное $2^{2x^3}$ и получаем $\frac1{2t}+32t-10=0$. После умножения на $2t$: $64t^2-20t+1=0$, $t=1/4$ или $1/16$. Из первого следует $x^2-3x+1=0$ и $x=(3\pm\sqrt5)/2$. Из второго $x^2-3x+2=0$ и $x=1$ либо $2$. На отрезке $[0{,}4;1{,}4]$ подходит только $1$: число $(3-\sqrt5)/2\approx0{,}382$ чуть меньше $0{,}4$, остальные больше $1{,}4$.}
\Heading{Проверка}{Все четыре найденных корня соответствуют положительным значениям $t$ и удовлетворяют исходному уравнению; проверка границ показывает, что внутри заданного отрезка находится ровно один.}
\Heading{Итог}{\Res{} Все корни: $\frac{3-\sqrt5}2$, $1$, $2$, $\frac{3+\sqrt5}2$; на $[0{,}4;1{,}4]$ только $1$.}
\Heading{Задача на закрепление}{Решите $2^{2x-1}+2^{3-2x}=5$ и укажите решения из отрезка $[0;2]$.}
\newpage
\phantomsection\label{task:3_16}
\section*{Лист 3, задание № 16. Логарифмическое неравенство}
\Heading{Текст задания}{Решите неравенство $\log_{\sqrt2}^{4}(24-x^2)-192\log_4^{2}(24-x^2)\ge64$.}
\Heading{Ваше решение}{Вы правильно начали с ОДЗ $24-x^2>0$, но при переходе к основанию $2$ заменили четвёртую степень первого логарифма множителем $2$, а второй коэффициент получили равным $96$, и решение осталось незаконченным.}
\Heading{Ваш ответ}{Ответ не указан}
\Heading{Ошибка}{\textcolor{red}{\textbf{Ошибка:}} Ошибочно возведён множитель перехода к новому основанию в степень.}
\Heading{Где именно возникает ошибка}{Условие $24-x^2>0$ определяет область допустимых значений и записано верно.}
\Heading{Первый неверный или недостающий шаг}{Первый неверный шаг --- преобразование $\bigl(2\log_2(24-x^2)\bigr)^4$ в $2\log_2^4(24-x^2)$: правильно $16\log_2^4(24-x^2)$.}
\Heading{Почему это неверно}{При смене основания $\log_{\sqrt2}z=2\log_2z$ и $\log_4z=\frac12\log_2z$. В четвёртой степени коэффициент $2$ превращается в $16$, а во второй степени коэффициент $1/2$ превращается в $1/4$. Нельзя возводить в степень только логарифм, оставляя множитель неизменным.}
\Heading{Ключевая идея}{\Idea{} Сначала задайте $L=\log_2(24-x^2)$ и отдельно вычислите коэффициенты $2^4$ и $(1/2)^2$.}
\Heading{Исправление от последнего правильного шага}{С учётом ОДЗ получаем $16L^4-192\cdot\frac14L^2\ge64$, то есть $L^4-3L^2-4\ge0$.}
\Heading{Полное правильное решение}{Положим $u=L^2\ge0$. Тогда $(u-4)(u+1)\ge0$, значит $u\ge4$ и $|L|\ge2$. Два случая: $L\ge2$ даёт $24-x^2\ge4$, следовательно, $|x|\le2\sqrt5$; $L\le-2$ даёт $0<24-x^2\le1/4$, то есть $95/4\le x^2<24$, откуда $\frac{\sqrt{95}}2\le|x|<2\sqrt6$. Объединяем три отрезка и полуинтервала.}
\Heading{Проверка}{При $x=0$: $L=\log_2 24>2$, неравенство выполняется. При $x=\pm2\sqrt6$ аргумент логарифма равен нулю, эти точки исключены. Границы $x=\pm2\sqrt5$ и $x=\pm\sqrt{95}/2$ включены по знаку нестрогого неравенства.}
\Heading{Итог}{\Res{} Решение: $[-2\sqrt5;2\sqrt5]\cup(-2\sqrt6;-\frac{\sqrt{95}}2]\cup[\frac{\sqrt{95}}2;2\sqrt6)$.}
\Heading{Задача на закрепление}{Решите $\log_{\sqrt2}^{4}(10-x^2)-48\log_4^{2}(10-x^2)\ge16$.}
\par\vfill
\begin{center}\small Лёвин Артём Александрович эксклюзивно для Николь Саркисьянц\end{center}
\end{document}
